\section{评估}
\label{sec:evaluation}

\subsection{测试设计}
\label{sec:test-design}

评估是一次小规模内部试点，不是用户研究。两名审阅者各自完成 48 份文档：一半先用
人工流程，一半先用助手流程，顺序随机化以减少学习效应。记录三项指标：中位审阅时间、
审阅者推翻模型提议的比例，以及证据完整度（审阅记录中能找到支持最终标签的原文片段
的比例）。所有文档均为合成数据，评估不涉及真实个人。

\subsection{结果}
\label{sec:results}

表~\ref{tab:results} 汇总结果。助手把中位审阅时间从 6.5 分钟降到 3.8 分钟，证据完整度
从 62\% 提升到 91\%。代价是 48 条中有 9 条被审阅者推翻，推翻率 18.8\%。这个数字不应
被读成「模型有 81.2\% 的准确率」：推翻发生在模型与审阅者意见不一致时，其中既有模型
错误，也有审阅者改变判断的情况。

\begin{table}[htbp]
  \centering
  \begin{tabular}{lccc}
    \toprule
    指标 & 人工流程 & 助手流程 & 变化 \\
    \midrule
    中位审阅时间（分钟） & 6.5 & 3.8 & $-42\%$ \\
    审阅者推翻率 & -- & 18.8\% & -- \\
    证据完整度 & 62\% & 91\% & $+29$ 个百分点 \\
    \bottomrule
  \end{tabular}
  \caption{小规模内部试点的结果。样本为 48 份合成文档，两名审阅者。}
  \label{tab:results}
\end{table}

\subsection{失败案例}
\label{sec:failures}

失败主要集中在三类情况：
\begin{itemize}
  \item \textbf{混合主题。} 文档同时涉及两个主题时，模型只给一个标签，审阅者需要拆分；
  \item \textbf{证据缺失。} 文档很短，模型给出高分但没有可引用的原文片段；
  \item \textbf{标签歧义。} 两个主题的边界本身不清楚，模型与审阅者的分歧无法归因于
        模型错误。
\end{itemize}

附录表~\ref{tab:test-cases} 给出逐条测试记录。我们保留这些失败案例，而不是只报告成功
样本，因为审阅者需要知道系统在哪些情况下不可靠。
